\section{Why the positive proofs stay inside iPRA}\label{sec:proof-audit}

The claim of an $\ipra$ proof is more specific than the claim that the
displayed numerical algorithms are primitive recursive. The arithmetic
must also prove their correctness with its permitted induction. We give
the reduction explicitly, so that an analytic-looking intermediate line
does not silently introduce stronger induction or a completeness axiom.

\subsection{Finite algebra certificates and approximation}
An integer polynomial identity in a fixed finite number of variables is
proved by expanding, collecting like monomials, and applying the
commutative-ring laws. These laws for encoded integers and rationals
follow by quantifier-free induction from the PR equations for addition
and multiplication. A finite list of coefficients is itself a PR object.
If one wants a uniform identity-checking procedure, polynomial addition,
multiplication and evaluation are PR finite-list operations; the proof
that evaluation respects them is a quantifier-free induction after the
finite lists are coded.

Every use of an orthogonality hypothesis in
Section~\ref{sec:universal} has an explicit polynomial residual identity.
For example, for an arbitrary matrix $A$, put
$E=A^TA-I$ and $D=\det A-1$. The exact identity
\begin{equation}\label{eq:adj-residual}
 \adj A-A^T=D A^T-E\adj A
\end{equation}
follows from $A\adj A=(\det A)I$ by multiplication with $A^T$.
Thus approximate orthogonality and determinant one imply approximate
cofactor inversion with a finite, computable error bound. The remaining
universal identities are obtained by the displayed additions and
multiplications, so the same argument applies to each one. There is no
step that converts determinant zero into a nonzero kernel vector by a
classical rank theorem.

Here is a general bound that makes this elimination of real notation
mechanical. If $P(X)=\sum_\alpha a_\alpha X^\alpha$, $B\ge1$,
$|x_i|,|y_i|\le B$, and $\max_i|x_i-y_i|\le\eta$, then
\begin{equation}\label{eq:poly-lipschitz}
 |P(x)-P(y)|\le L(P,B)\eta,\qquad
 L(P,B)=\sum_\alpha |a_\alpha|\,|\alpha|\,B^{|\alpha|-1},
\end{equation}
where constant monomials contribute zero. To prove this, write the
difference of a product of $m$ factors as a sum of $m$ products with one
difference factor. Each such term is bounded by $B^{m-1}\eta$.
Sum over the monomials. This is finite telescoping and finite summation.
Rational coefficients cause no difficulty after positive denominators
are cleared.

For products, quotients with a displayed lower bound, and square roots,
use Lemma~\ref{lem:operations}. Reading an expression tree from its root
to its inputs assigns an explicit required input precision at each
node. A product multiplies its error budget by the displayed magnitude
bound; a quotient uses \eqref{eq:division-error}; a square root requests
twice the precision plus a fixed shift. All expression trees in the
matrix proof have fixed finite size. Their resulting precision
transformers are PR in $n,h$ and the finite magnitude bounds.

For instance, after clipping the nine rotation-entry approximants to
$[-1,1]$, entry error $\eta$ gives error at most $3\eta$ in every entry
of $S$, at most $2\eta$ in $d$, and at most $\eta$ in every component
of $w$. With $\delta=2^{-h}$, sample the input at precision $h+6$.
Then $3\eta\le3\delta/64<\delta/12$, as required by the pivot proof.
Clipping is a rational min/max operation and does not increase error
when the represented entry lies in $[-1,1]$. Orthogonality proves that
entry bound by a sum of squares.

The chosen column has $q=v\cdot v=dp\ge\delta^2/6$ and
$\sqrt q\ge\delta/3$. Thus every later precision is bounded before
execution. The stored index $j$, obtained by three rational comparisons,
is independent of the requested output precision. This proves the
regularity of the output name as well as its algebraic properties.

\subsection{Where induction occurs}
The following table lists the induction tasks. ``Finite'' means a
primitive-recursive bound is part of the construction.

\begin{longtable}{@{}p{0.27\textwidth}p{0.64\textwidth}@{}}
\toprule
Task & Quantifier-free invariant or identity \\
\midrule
\endhead
Integer and rational arithmetic & Defining equations for addition,
multiplication and comparison; denominator positivity and the cleared
integer equality or inequality.\\[3pt]
Finite sums and products & Equality between the accumulator and the
finite folded expression; triangle or product error bounded by the
corresponding finite error sum.\\[3pt]
Maximum and pivot & The selected rational sample dominates every sample
already processed. There are only three axis samples.\\[3pt]
Integer square root & Current candidate is in range and has admissible
square; every larger processed candidate is inadmissible, encoded by
a bounded Boolean fold.\\[3pt]
Taylor estimates & Factorial-ratio bounds, the finite geometric-sum
bound, and the finite binomial convolution at the requested degree.\\[3pt]
Arctangent mesh comparison & The finite rational difference between
two adjacent dyadic meshes has the displayed bound.\\[3pt]
Cayley product & The difference of the two products of the first $j$
factors is bounded by the sum of their first $j$ errors.\\[3pt]
PR input simulation in the obstruction & Deterministic bounded
simulation has at most one first halting step and one output.\\
\bottomrule
\end{longtable}

When an invariant contains a bounded ``for every processed index'',
replace it with its PR Boolean fold. The induction formula is then
quantifier free in the full PR language. Its context may contain
universally quantified assumptions about the input names; the
quantifier-free restriction is on the induction formula, not on the
whole context.

For additional precision, fix a requested output precision $n$ before
performing an induction. The finitely many input samples needed by that
calculation form a coded finite table. Its regularity and polynomial
residual bounds, at the sampled indices, are decidable rational
conditions. Let $\mathsf{GoodTable}$ be their Boolean conjunction. The
actual induction can be applied to a formula of the form
\[
 \mathsf{GoodTable}=1\ \longrightarrow\
 \mathsf{AccumulatorBound}(j,n,h,\mathsf{table})=1.
\]
All quantities and loop bounds are PR. Instantiate the input's
universally quantified premises at those indices to establish
$\mathsf{GoodTable}=1$. Only after finishing the finite calculation do
we universally quantify over $n$. We never apply induction to a
statement asserting a Cauchy property at all precisions.

\subsection{The actual arithmetic conclusions}
For rational input the first endpoint is particularly simple. Encode a
rational matrix by a natural number. Let $\mathsf{RatSO3}(r)$ be the
decidable test of its nine orthogonality equations and its determinant.
Let $V(r),Q(r),C(r),D(r),B(r)$ be the totalized PR extraction functions.
The statement
\[
 \forall r\bigl(\mathsf{RatSO3}(r)\longrightarrow
   \mathsf{RatCertificate}(r,V(r),Q(r),C(r),D(r),B(r))\bigr)
\]
is an ordinary arithmetic theorem, with a quantifier-free matrix of
decoded rational equalities and inequalities. The predicate
$\mathsf{RatCertificate}$ is exactly \eqref{eq:ratcert}, together with
$Q=V\cdot V>0$ and $D=1-C$. The normalization program adds universal
finite-precision bounds; interpreting its explicit regular names gives
the exact unit-axis assertion.

For each fixed PR real-matrix family $A$, the separated result has the
form
\begin{equation}\label{eq:object-schema}
\begin{split}
 \forall\mathbf p\,\forall h\bigl(&\mathsf{Regular}(A,\mathbf p)
 \wedge\mathsf{Orthogonal}(A,\mathbf p)
 \wedge\mathsf{DetOne}(A,\mathbf p)
 \wedge\mathsf{Gap}(A,\mathbf p,h)\\
 &\longrightarrow\mathsf{UnitAxisAndRodrigues}
       (A,U_A,C_A,S_A,\mathbf p,h)\bigr).
\end{split}
\end{equation}
Each displayed predicate abbreviates its first-order rational
approximation formula; its unbounded quantifiers range over numerical
precisions. The output symbols are the concrete PR definitions already
constructed. There is no existential quantifier ranging over an
unspecified function space. The finite proof described above is a
derivation of each such schema instance. The unconditional approximation
and quaternion statements have the same form, with the gap replaced by
the robust finite branch supplying it where needed.

This explains both parts of the logical status. The positive proofs
translate into the stipulated arithmetic because every computational
loop and invariant is explicit. The unrestricted exact unit-axis claim
does not become true by calling it a schema: the fixed PR family in
Section~\ref{sec:obstruction} is already a forbidden instance.

\section{Lean formalization: existing interfaces and a completion path}
\label{sec:lean}

\subsection{Pinned repository and reusable infrastructure}
The repository was inspected at commit
\begin{center}\small\ttfamily
c1bda76f7f86d2ade505355980826f625e36fff3.
\end{center}
The root toolchain and Mathlib dependency both specify version
\texttt{v4.32.0}. The relevant source names and their roles are:
\begin{center}
\begin{tabular}{@{}p{0.40\textwidth}p{0.52\textwidth}@{}}
\toprule
Interface & Role \\
\midrule
\texttt{PR}, \texttt{PTm} & PR function symbols and first-order terms.\\
\texttt{Fm}, \texttt{Prf} & Formulas and intuitionistic derivations.\\
\texttt{PRAAx.ind}, \texttt{PRA.qfInd} & Induction with an explicit
\texttt{QF} side condition.\\
\texttt{pra\_theorem\_iff} & Transport from iPRA to the scalar theory
$T_{\rm eq}$.\\
\texttt{pra\_iff\_native} & Transport from iPRA to the native concise
calculus.\\
\bottomrule
\end{tabular}
\end{center}
These are actual inspected declarations \cite{proveit}. They transport
an existing object-language proof. They do not translate an arbitrary
theorem about Mathlib's real numbers into an iPRA derivation.

\subsection{Artifacts supplied now}
\texttt{lean/RodriguesKernel.lean} contains substantive algebraic proof
scripts over an arbitrary commutative ring. They cover dot and cross
identities, the vector triple product, Lagrange's identity, fixation of
the unit axis, preservation of dot products, and determinant one. The
scripts use explicit coordinates and \texttt{ring}. They contain no
\texttt{sorry} declaration or added axiom, but they have \emph{not} been
run through Lean in this environment.

\texttt{lean/ObjectLanguageBoundary.lean} uses the repository's actual
\texttt{PR} type to specify rational-name symbols. It also gives a small
transport lemma taking an already supplied iPRA derivation to the
native calculus. This is explicitly an interface example, not the
missing geometry derivation.

The environment had no \texttt{lean} or \texttt{lake} executable.
Accordingly no claim is made that these source files elaborate, that
their dependencies have passed a kernel audit, or that the entire
theorem is formalized. The supplied detailed plan identifies the
remaining work; this is the requested outline with a concrete path to
completion.

\subsection{A clear sequence of formalization milestones}
\begin{enumerate}
\item \textbf{Compile the finite algebra prototype.} In the pinned
ProveIt environment, compile \texttt{RodriguesKernel.lean}; resolve any
syntax or simplifier differences. Inspect \texttt{\#print axioms} for
the final declarations. This checks the algebra specification but does
not yet establish object-language derivability.
\item \textbf{Implement encoded rational arithmetic.} Construct
\texttt{PR} definitions for signed integers, positive-denominator
rationals, comparison, powers, finite folds and bounded search. Prove
their arithmetic and order properties as \texttt{Prf PRAAx} terms,
using \texttt{PRA.qfInd} for the displayed finite invariants.
\item \textbf{Turn fixed polynomial identities into derivations.}
Either write direct ring-law derivations for the finite collection here,
or implement a certificate elaborator for coefficient normalization.
Its output must be an actual object-language derivation. A Python
symbolic success message and a semantic Lean \texttt{ring} proof are
independent checks, not substitutes for that output.
\item \textbf{Finish the rational theorem first.} Its input predicate,
identity test, pivot, and certificate are all finite and rational.
Prove the totalized extractor correct, then prove the dyadic-root
normalization interval and the precision bounds of
Section~\ref{sec:rational}. This gives a compact exact initial target.
\item \textbf{Formalize regular-name arithmetic.} Encode the predicates
of Section~\ref{sec:arithmetic} and prove the explicit operation
transformers. Instantiate universal input premises at the finitely
needed precision indices. Keep target precision free throughout every
induction, and universally introduce it afterward.
\item \textbf{Formalize the gap construction.} Prove the universal
cofactor identities, the three-sample pivot bound, and the selected
column's positive norm. Compose the name operations to obtain
$U_A,C_A,S_A$, then build \eqref{eq:object-schema}. The rank-one
argument is the vanishing-minor proof given above; no rank API or
eigenvector choice is needed.
\item \textbf{Add the robust branches.} Formalize the small/large
$d$ comparison for finite-precision Euler and the two $c$ charts for
the exact quaternion theorem. Their decisions are rational and their
proofs use the displayed overlapping thresholds.
\item \textbf{Add radian angles.} Define finite Taylor polynomials and
dyadic arctangent sums as PR functions. Prove (T1)--(T9) at fixed
precision using the listed invariants; add the two circle charts. This
proves both the given-angle version of Rodrigues and reconstruction
of a radian angle from the extracted circle pair.
\item \textbf{Transport only the finished derivation.} Close the
arithmetic schema instance, construct its actual
\texttt{[] $\vdash$[PRAAx] formula} proof term, and apply
\texttt{pra\_theorem\_iff} or \texttt{pra\_iff\_native}. Audit the
induction side conditions and ensure there is no new arithmetic axiom.
\end{enumerate}

The nonprovability theorem is a separate metamathematical development:
formalize bounded universal simulation and its unique first halt,
the Cayley reduction, and the realizability soundness theorem. It is
not needed to validate any positive algorithm. Its explicit arithmetic
target is \eqref{eq:H}; the note does not pretend that the concise
repository already contains its formal nonprovability proof.

\section{Verification, scope of evidence, and reproducibility}
\label{sec:verification}

The executable certificate module uses exact rational arithmetic.
It accepts rational $R$, verifies \eqref{eq:so}, extracts
$(v,q,c,d,b)$, and checks every equality in \eqref{eq:ratcert} exactly.
The normalization routine returns rational lower and upper square-root
bounds and a rational coordinate error certificate. It never treats
a floating-point tolerance as exact equality.

The test suite includes every nonzero integer quaternion in
$\{-2,-1,0,1,2\}^4$: 624 input quaternions, including four identity
inputs and 124 half-turn inputs. There are 31 additional edge inputs,
including small-angle and near-half-turn scales through $2^{128}$.
The initial normalization suite checks 1,965 combinations of those
inputs with requested precisions $0,8,32$, and rejects three invalid
matrices. The JSON reports record the final exact counts, including
the additional finite-precision branch checks.

The symbolic supplement checks the finite adjugate, dot/cross,
Rodrigues, Cayley and quaternion numerator identities by exact polynomial
arithmetic. These checks provide useful independent evidence for the
finite calculations. They do not prove the quantified $\ipra$ theorem,
certify the analytic estimates automatically, or establish Lean
elaboration. Those limits are stated in the source files and the
formalization status file.

To reproduce the computational evidence and PDF, from the package root:
\begin{lstlisting}
python3 python/verify.py
python3 python/verify_polynomials.py
pdflatex -interaction=nonstopmode -halt-on-error euler_rodrigues_ipra.tex
pdflatex -interaction=nonstopmode -halt-on-error euler_rodrigues_ipra.tex
\end{lstlisting}
The symbolic checker needs SymPy; the rational certificate checker uses
only Python's standard library. The package includes the LaTeX sources,
compiled PDF, Lean prototype and plan, exact-arithmetic scripts,
verification reports, and a source manifest. The manifest pins the
inspected repository files and the two relevant drafts. The repository
was read, not modified.

\begin{thebibliography}{9}
\raggedright
\small
\bibitem{eulerwiki}
Wikipedia contributors, \emph{Euler's rotation theorem}, user-specified
statement, accessed November 14, 2026.
\url{https://en.wikipedia.org/wiki/Euler%27s_rotation_theorem}.

\bibitem{rodrigueswiki}
Wikipedia contributors, \emph{Rodrigues' rotation formula}, user-specified
statement, accessed November 14, 2026.
\url{https://en.wikipedia.org/wiki/Rodrigues%27_rotation_formula}.

\bibitem{proveit}
Vladimir Reshetnikov, \emph{ProveIt: ConciseIPRA}, repository snapshot
\texttt{c1bda76f7f86d2ade505355980826f625e36fff3}.
Inspected \texttt{README.md}, \texttt{PR.lean}, \texttt{Syntax.lean},
\texttt{BiInterpretation.lean}, \texttt{NativeBiInterpretation.lean},
and the root toolchain/build configuration.
\url{https://github.com/VladimirReshetnikov/ProveIt/tree/c1bda76f7f86d2ade505355980826f625e36fff3/Logic/PrimitiveRecursiveArithmetic/ConciseIPRA}.

\bibitem{drafts}
ProveIt research drafts, \texttt{docs/drafts/FTA\_iPRA.zip} and
\texttt{docs/drafts/triangle\_formulas\_ipra.zip}, at the same pinned
snapshot. Inspected their source discussions of PR-name schemas,
quantifier-free induction and formalization status. No geometric
result from these drafts is assumed here.
\url{https://github.com/VladimirReshetnikov/ProveIt/tree/c1bda76f7f86d2ade505355980826f625e36fff3/docs/drafts}.

\bibitem{bridgesrichman}
Douglas Bridges and Fred Richman, \emph{Varieties of Constructive
Mathematics}, London Mathematical Society Lecture Note Series 97,
Cambridge University Press, 1987, Chapter 1 (LLPO and constructive
existence). Publisher's excerpt:
\url{https://assets.cambridge.org/97805213/18020/excerpt/9780521318020_excerpt.pdf}.

\bibitem{kleene}
S.~C.~Kleene, ``On the interpretation of intuitionistic number theory,''
\emph{Journal of Symbolic Logic} \textbf{10}(4) (1945), 109--124.
\url{https://doi.org/10.2307/2269016}.
Author's article in an academic course archive:
\url{https://webspace.science.uu.nl/~ooste110/studsemIntMod/Kleene45.pdf}.

\bibitem{moschovakis}
Joan Rand Moschovakis and Garyfallia Vafeiadou, \emph{Intuitionistic
Mathematics and Logic}, authors' English translation of their 2009
chapter, section on number realizability and Nelson's theorem.
\url{https://nuprl-web.cs.cornell.edu/documents/Rand/intuitionistic-math2.pdf}.
\end{thebibliography}
